% ZeroDDS Architecture — Reference.
%
% Build: tectonic architecture.tex   →  architecture.pdf
% Or via top-level Makefile:          make -C documentation pdf-arch

\documentclass[11pt,a4paper]{article}
\usepackage[margin=2.4cm]{geometry}
\usepackage{fontspec}
\usepackage{microtype}
\usepackage{xcolor}
\usepackage{listings}
\usepackage{hyperref}
\usepackage{titlesec}
\usepackage{enumitem}
\usepackage{tabularx}
\usepackage{booktabs}
\usepackage{longtable}
\usepackage{tikz}

\setmainfont{Source Serif 4}[
  UprightFont = *,
  BoldFont    = *-Bold,
  ItalicFont  = *-Italic,
]
\setsansfont{Source Sans 3}
\setmonofont{Source Code Pro}[Scale=MatchLowercase]

\definecolor{accent}{HTML}{1F3A8A}
\definecolor{codebg}{HTML}{F4F4F4}

\hypersetup{
  colorlinks = true,
  linkcolor  = accent,
  urlcolor   = accent,
  pdfauthor  = {ZeroDDS Maintainers},
  pdftitle   = {ZeroDDS Architecture Reference},
  pdfsubject = {ZeroDDS doc-trail station 02 — Architecture},
}

\titleformat*{\section}{\Large\bfseries\sffamily\color{accent}}
\titleformat*{\subsection}{\large\bfseries\sffamily}
\titleformat*{\subsubsection}{\normalsize\bfseries\sffamily}

\lstset{
  basicstyle    = \small\ttfamily,
  backgroundcolor = \color{codebg},
  frame         = single,
  rulecolor     = \color{codebg},
  framerule     = 0.4pt,
  framesep      = 4pt,
  breaklines    = true,
  showstringspaces = false,
  columns       = fullflexible,
}

\title{\sffamily\bfseries\color{accent}\Huge ZeroDDS\\\Large Architecture Reference}
\author{\sffamily ZeroDDS Maintainers}
\date{\sffamily \today}

\begin{document}

\maketitle
\vspace{1em}

\begin{quote}\itshape
This is the print companion to documentation/02-architecture/.
The Markdown source is the canonical version; this PDF is rendered
on demand via tectonic for offline review and printed handout.
\end{quote}

\tableofcontents
\newpage

\section{Mission}

ZeroDDS is a pure-Rust implementation of OMG DDS 1.4, DDSI-RTPS 2.5,
DDS-Security 1.2, and DDS-XTypes 1.3, with native bindings for
C, C++, C\#, Java, Python, TypeScript, and a ROS-2 RMW plugin.

The name encodes the contract: \emph{zero dependencies} (Safe-Core),
\emph{zero panic} (contract), \emph{zero unsafe} (where structurally
possible), \emph{zero copy} (SHM path), \emph{zero vendor lock-in}.

\section{Layered architecture}

\begin{lstlisting}
+-----------------------------------------------------------+
| Bindings        C | C++ | C# | Java | Python | TypeScript |
+-----------------------------------------------------------+
| DCPS API        dcps                                      |
| Bridges         coap | websocket | mqtt | amqp | grpc |   |
|                 ros2-rmw | opcua | soap                   |
+-----------------------------------------------------------+
| Discovery       discovery (SPDP / SEDP / WLP)             |
| Wire            rtps (RTPS 2.5)                           |
| Security        security-{pki, crypto, rtps, ...}         |
+-----------------------------------------------------------+
| Encoding        cdr | idl | types | qos                   |
+-----------------------------------------------------------+
| Transport       transport-{udp, tcp, shm, uds, tsn}       |
+-----------------------------------------------------------+
| Foundation      buffer | crc | rcu | observability |      |
|                 rt-linux                                  |
+-----------------------------------------------------------+
\end{lstlisting}

A request flows top-down on the publisher side and bottom-up on
the subscriber side. The hot path is
\texttt{dcps $\rightarrow$ rtps $\rightarrow$ transport-udp};
everything else (bridges, security, observability) is opt-in
indirection.

\section{Process model}

A single \texttt{DcpsRuntime} owns:

\begin{itemize}[leftmargin=*]
  \item One UDP socket pair per discovery channel: SPDP multicast
        \texttt{239.255.0.1:7400 + 250 $\times$ domain\_id},
        SPDP unicast, and a separate user-data unicast socket.
  \item An async event-loop thread that pumps inbound datagrams,
        handles SPDP/SEDP, and ticks the reliable writer/reader
        state machines at \texttt{tick\_period} (default 50 ms).
  \item A registry of user writers and readers using
        \texttt{Arc<RwLock<BTreeMap<EntityId, Arc<Mutex<Slot>>>>>} so
        that two app threads writing to different topics never
        serialise on a global lock.
\end{itemize}

User code calls \texttt{register\_user\_writer} /
\texttt{register\_user\_reader} to attach endpoints, and
\texttt{write\_user\_sample} / receives via an
\texttt{mpsc::Receiver} to push and pull samples.

\section{Wire stack}

Every byte that leaves a participant traverses the same stack:

\begin{enumerate}[leftmargin=*]
  \item Sample bytes (CDR / XCDR2-encoded by the IDL stubs).
  \item User-payload encapsulation: a 4-byte header
        \texttt{00 07 00 00} = XCDR2 little-endian, options=0.
  \item DATA or DATA\_FRAG submessage carrying the encapsulated
        bytes, plus writer EID, reader EID, and SN.
  \item Larger samples are fragmented; each fragment lives in its
        own DATA\_FRAG submessage.
  \item RTPS message header (16 bytes): \texttt{"RTPS"} magic,
        version 2.5, vendor ID, GUID prefix.
  \item Optional: SRTPS wrap (header AAD + encrypted body) per
        DDS-Security 1.2 §7.4.6.6.
  \item UDP datagram, unicast or multicast.
\end{enumerate}

The exact wire format is DDSI-RTPS 2.5; ZeroDDS is byte-identical
to Cyclone DDS, FastDDS, RTI Connext Pro, and OpenSplice.

\section{Hot path: \texttt{write\_user\_sample}}

The hot path on the publisher side, with the locks taken at each
step:

\begin{lstlisting}
DcpsRuntime::write_user_sample
  1. read-lock user_writers RwLock      (registry)
  2. clone Arc<Mutex<Slot>>; drop read   (no global hold)
  3. lock per-slot Mutex                 (per-writer)
  4. branch on size:
     - small (<= 1.5 KiB): PoolBuffer<1536> on the stack
     - big: Vec::with_capacity fallback
  5. ReliableWriter::write
       - SN++
       - Arc::from(&[u8])  (the only heap alloc on the path)
       - cache.insert(CacheChange::alive_arc(SN, payload))
       - cache_stats.refresh()           (atomic counters)
       - per-proxy DATA submessage build
  6. drop slot Mutex
  7. per-target security wrap (if feature on)
  8. send_on_best_interface
\end{lstlisting}

Key invariants:

\begin{itemize}[leftmargin=*]
  \item Two locks total on the hot path; both are bounded.
  \item Per-slot \texttt{Mutex} on every user-writer/reader — two
        writers on different topics never serialise.
  \item Atomic stats on the history cache — monitoring threads can
        poll latency / fill-level without disturbing the writer.
  \item For small samples, only one heap allocation
        (\texttt{Arc::from}). Encap framing avoids heap touches
        via \texttt{PoolBuffer}.
\end{itemize}

\section{Concurrency model}

\begin{itemize}[leftmargin=*]
  \item Per-slot \texttt{Mutex} on every user-writer and
        user-reader.
  \item Atomic stats on the history cache —
        monitoring threads poll \texttt{Arc<HistoryCacheStats>}
        without taking any lock.
  \item Optional lock-free read cache — opt-in
        \texttt{LockFreeReadHistoryCache} for read-heavy paths via
        copy-on-write \texttt{Arc<Inner>} swap.
\end{itemize}

\section{Real-time profile}

\texttt{crates/rt-linux} (UNSAFE-FFI, isolated to a single crate
with explicit \texttt{// SAFETY:} commentary on every unsafe
block) wraps Linux \texttt{sched\_setattr(2)} for SCHED\_FIFO,
SCHED\_RR, SCHED\_DEADLINE, plus \texttt{sched\_setaffinity(2)}
for CPU pinning. The accompanying deployment guide
(\texttt{docs/REALTIME\_DEPLOYMENT.md}) covers
\texttt{isolcpus}, \texttt{nohz\_full}, \texttt{rcu\_nocbs},
preempt\_rt, and a per-CPU layout suggestion for an 8-core host.

\section{What's not in this PDF}

\begin{itemize}[leftmargin=*]
  \item Per-crate API reference — see the rustdoc bundled at
        \texttt{documentation/api/} or
        \texttt{cargo doc --workspace}.
  \item Internal-developer architecture (German) — see
        \texttt{docs/architecture/00\_overview.md} through
        \texttt{09\_delegation.md}.
  \item Configuration reference — see trail station 03.
  \item IDL reference — see trail station 04.
  \item Per-language integration — see trail station 05.
\end{itemize}

\vfill
\begin{center}
\itshape ZeroDDS Documentation Trail — station 02 (architecture).\\
The Markdown source is canonical; this PDF is rendered from it
on demand.
\end{center}

\end{document}
